\chapter{Diskussion}
\section{Bewertung der Datenqualität}
Hier sollen Außergewöhnlichkeiten in den Daten sowie Umstände in der Auswertungsmethodik diskutiert werden. Diese können Einfluss auf die Qualität der Messergebnisse haben.\newline

Bei der Auswertung des P1-Paradigmas wird der Mittelwert der Spannung im Bereich 80 - 200 ms nach dem Reiz zwischen den Systemen verglichen. Wie in Abbildung \ref{erps_ergebnisse} zu sehen, fällt die Spannung im unteren Diagramm für das BioSemi-System am Ende des betrachteten Zeitraums schneller ab als die Spannung der OpenBCI-Systeme. Das hat zur Folge, dass bei der Mittlung der Wert für das BioSemi-System geringer ist als die tatsächliche durchschnittliche Amplitude der P1. 
Dieser Umstand verfälscht die Daten dahingehend, dass der Wert der Spannung bei kleinen Bildern für BioSemi größer ist als . Demnach ist das Verhältnis s1/s3 und s2/s3 geringer.\newline
Bei den großen, kontrastreichen Bildern ist der Unterschied zwischen der Negativität der Kurven kurz vor dem Ende des Intervalls sehr gering, hat also keine Auswirkung auf die Ergebnisse.\newline

Eine weitere Außergeöhnlichkeit fällt bei der Frequenzbandanalyse auf.
In beiden OpenBCI-Durchläufen tritt ein schmaler Peak, also ein Hochpunkt, bei 25 Hz auf, der beim BioSemi-System fehlte (in Abbildung \ref{frequenzbandanalyse} zu sehen). Er ist zu groß, zu regelmäßig und unabhängig von Augenzustand und Bedingung (keine Unterschiede in Amplitude oder Frequenz zwischen \acrshort{eo}/\acrshort{ec} sowie zwischen s1 und s2), um neuronalen Ursprungs zu sein. Auch ist nach aktuellem Stand kein physiologischer Effekt bekannt, der einen solchen Ausschlag erklären würde. Für eine technische Ursache kommen mehrere Möglichkeiten infrage:\newline
\begin{itemize}
  \item Netzbrummen scheidet aus, da der Peak nicht bei 50 Hz liegt und sich in der Kabine keine Netzteile befinden. Gegenüber dem BioSemi-Durchlauf kam lediglich der Laptop hinzu.
  \item Bildschirmflackern: Bei Pixel-Inversion oder Frame Rate Control entsteht ein Flackern mit der halben Bildwiederholrate, bei einem 60-Hz-Monitor also bei 30 Hz. %(Quelle: flickersense.org/background/led-screens). Das erklärt einen Peak bei 25 Hz nicht direkt.
  \item Das Cyton-Board könnte aus anderen Gründen ein elektromagnetisches Feld mit 25 Hz abstrahlen.
\end{itemize}

Mit Sicherheit lässt sich die Ursache auf Basis der vorliegenden Daten nicht bestimmen. Am wahrscheinlichsten ist ein Einfluss des Boards, möglicherweise auch des Laptops. Der Peak hat keinen Einfluss auf die Auswertung der Daten. Bei der Programmierung des Roboterarms ist jedoch zu beachten, dass diese Frequenz von 25 Hz nicht für die Bewegung des Armes genutzt werden sollte.\newline

\section{Bezug auf die Hypothesen}

Das ist ein klarer Beleg für den \gls{berger-effekt}. Zudem spricht die räumliche Verteilung gegen einen äußeren Störeinfluss, da ein solcher alle Cluster gleich stark betreffen würde. Stattdessen war die Suppression systematisch dort am größten, wo der Berger-Effekt physiologisch zu erwarten ist.\newline

Die Alpha-Suppression war mit dem OpenBCI-System bei allen Versuchspersonen deutlich ausgeprägt, auch bei offener Kabinentür. Das spricht dafür, dass sich das Öffnen und Schließen der Augen grundsätzlich als einfaches Steuersignal für einen Roboterarm eignen könnte. Ob dies auch in Echtzeit zuverlässig gelingt, muss in einer Folgestudie geprüft werden.



\section{Vor- und Nachteile eines tragbaren EEG-Systems}
\section{Vor- und Nachteile eines stationären EEG-Systems}